% Prepared for Edward Baker, 4 November 2026.
% Substantial LLM assistance: gpt-6-astra, actual effort ultra.
% Optional manuscript section, based on frozen RGD Note14.
% Accepted within scope by separate internal RGD-E04 r7 assessment; no specialist review.

\section{Endpoint-preserving pair comparison and a diagonal obstruction}
\label{sec:pair-comparison}

For a generating function $F(z)=\sum_{j\ge0}f_jz^j$, with $f_0=1$,
write
\[
 \Delta_{r,k}(F)=\det[f_{k+j-i}]_{0\le i,j<r},\qquad f_j=0\quad(j<0).
\]
The original normalization is $D_{r,k}=c_0^r\Delta_{r,k}(E)$.
The comparison below preserves $E(0)$ and $E(1)$ and has the correct
entire-function growth. The subsequent obstruction concerns only the
listed algebraic data; it is not a negative determinant for the theta
source.

\begin{proposition}[Endpoint-preserving pair comparison]
\label{prop:endpoint-pair}
Suppose
\[
 E(z)=\prod_i(1+a_i z)\prod_\ell
       (1+2s_\ell\cos\theta_\ell\,z+s_\ell^2z^2),
\]
where $a_i,s_\ell>0$, $0<\theta_\ell<\pi/2$, the complete alphabet is
infinite, $\sum_i a_i+\sum_\ell s_\ell<\infty$, and
$\sum_\ell\tan^2\theta_\ell<\infty$.
Set
\[
 y_\ell=\sqrt{1+2s_\ell\cos\theta_\ell+s_\ell^2}-1,
 \quad d_\ell=s_\ell^2-y_\ell^2,
 \quad
 R(z)=\prod_i(1+a_i z)\prod_\ell(1+y_\ell z)^2.
\]
Assume also $y_\ell\le2$. Then $R(0)=E(0)=1$, $R(1)=E(1)$, and,
with
\[
 V=\prod_\ell\frac{s_\ell^2}{y_\ell^2}<\infty,
\]
one has
\[
 1\le\frac{E(t)}{R(t)}\le V\quad(t\ge1),
 \qquad \lim_{t\to\infty}\frac{E(t)}{R(t)}=V.
\]
If $E-R=\sum_j\varepsilon_jz^j$, then
\begin{equation}
 \sum_{j\ge0}|\varepsilon_j|t^j\le(V-1)R(t)
 \qquad(t\ge0).
 \label{eq:pair-coeff-payment}
\end{equation}
Writing $(b_j)$ for the positive alphabet of $R$, every rectangular
reference determinant satisfies
\begin{equation}
 \Delta_{r,k}(R)\ge(b_1\cdots b_k)^r>0.
 \label{eq:pair-reference-margin}
\end{equation}
Let $\Theta=\sup_\ell\theta_\ell$, with $\Theta=0$ if there are no
nonreal pairs. If $(r+1)\Theta<\pi$, put
\[
 \kappa_r=\frac{\sin((r+1)\Theta)}{(r+1)\Theta},
 \qquad \kappa_r=1\text{ when }\Theta=0.
\]
Then
\begin{equation}
 \kappa_r^{rk}\le
 \frac{\Delta_{r,k}(E)}{\Delta_{r,k}(R)}
 \le\sec(\Theta)^{rk}.
 \label{eq:pair-sector-margin}
\end{equation}
\end{proposition}

\begin{proof}
For one pair, abbreviate its quadratic factor by $F$ and $(1+yz)^2$
by $G$. Direct algebra gives
\[
 s\cos\theta\le y<s,\qquad
 F-G=d\,z(z-1),\qquad 0<d/y^2\le\tan^2\theta.
\]
Thus the factors agree at zero and one. All products converge locally
uniformly. For $t\ge1$, $F(t)/G(t)$ increases from one to $s^2/y^2$.
Moreover
\[
 0\le2\log(s/y)\le-2\log\cos\theta\le\tan^2\theta.
\]
The summability assumption proves the finite quotient bound and its
limit by dominated convergence of the logarithms.

For finite products, expand the exact difference using
$F=G+d\,z(z-1)$. After taking coefficient absolute values, replace each
selected $z(z-1)$ by $z(z+1)$. At $t\ge0$ this gives the upper bound
\[
 R(t)\left\{\prod_\ell
 \left(1+\frac{d_\ell t(t+1)}{(1+y_\ell t)^2}\right)-1\right\}.
\]
Since $y\le2$ implies $y^2t(t+1)\le(1+yt)^2$, this is at most
$(V-1)R(t)$. Coefficient convergence and the same finite bound justify
passage to the infinite product.

The Toeplitz matrix of a finite positive-alphabet product is a product
of upper bidiagonal matrices $I+b_jS$. Their minors are nonnegative.
For the rectangle with row indices $0,\ldots,r-1$ and column indices
$k,\ldots,k+r-1$, Cauchy--Binet includes the positive term obtained by
shifting the whole block once in each of the first $k$ factors. Its
value is $(b_1\cdots b_k)^r$. Taking the infinite-product limit proves
\eqref{eq:pair-reference-margin}, including all negative-index zeros.

For completeness, rectangular Jacobi--Trudi identifies
$\Delta_{r,k}(E)$ with $s_{(r^k)}$ evaluated on its complete alphabet.
Classical Schur branching~\cite{EtingofRepresentationTheory20110110} into the real singleton and conjugate-pair
alphabets has nonnegative integer multiplicities. For a two-variable
partition $\nu=(u,v)$, $u\ge v\ge0$, the exact ratio is
\[
 \frac{s_\nu(se^{i\theta},se^{-i\theta})}{s_\nu(y,y)}
 =\left(\frac{s}{y}\right)^{u+v}
   \frac{\sin((u-v+1)\theta)}{(u-v+1)\sin\theta}.
\]
Each component appearing in the branching has $u\le r$: every
individual variable in $s_{(r^k)}$ has degree at most $r$, and the
nonnegative character multiplicities cannot cancel a highest-weight
monomial. Thus $u-v+1\le r+1$. Under the angular hypothesis, sinc
monotonicity, $\sin\theta\le\theta$, and $s/y\le\sec\Theta$ bound
each nonempty pair ratio below by $\kappa_r$ and above by
$\sec(\Theta)^{u+v}$. For the upper bound one also uses
$|\sin(j\theta)|\le j\sin\theta$, which follows from the finite
geometric sum. There are at most $rk$ nonempty pair components, and
their total degree is at most $rk$. Dividing the branching expansion
by its positive reference value expresses the full ratio as a
nonnegative weighted average of these products. This proves
\eqref{eq:pair-sector-margin}; local convergence gives the infinite
version.
\end{proof}

For the normalized theta function, the genus-zero product~\cite{Katkova0505174v1} is indexed
by upper-half-plane zeros $\rho=\sigma+i\gamma$ with multiplicity,
using $x=\alpha^{-1}$ and $\alpha=-(2\rho-1)^2$. The usual strip and
zero-counting inputs~\cite{Odlyzko1987AuthorCopy}, together with the inherited fact $\gamma>14$,
give $\Re x>0$, $|x|<1/784$, and
\[
 |\arg\alpha|=2\arctan\frac{|2\sigma-1|}{2\gamma},
 \qquad\sum\tan^2(\arg\alpha)<\infty.
\]
They verify the proposition's assumptions without RH or simplicity.
Both generating functions have nonnegative coefficients, so the
bounded positive-axis quotient also matches their maximum-modulus
growth up to a bounded factor. The reference nevertheless uses the
complete, implicit zero data; it is not an acquired numerical model.
No positive cosh-transform source or positive Green primitive for this
reference is asserted. The bounded-order estimate is classical sector
positivity, not a new theta positivity range. Requiring all pair Schur
blocks to be nonnegative at every degree would force every angle to
vanish and would restate RH for this product.

\begin{theorem}[Small complete scalar defects can have negative diagonal minors]
\label{thm:diagonal-pair-control}
Fix $B\in(1,2)$. For every integer $q\ge26$ there exist real entire
functions $E_q$ and $R_q$, with positive Taylor coefficients, such that:
\begin{enumerate}
 \item $E_q(0)=R_q(0)=1$ and $E_q(1)=R_q(1)=B$;
 \item both functions have genus zero, order $1/2$, and radial zero
 count asymptotic to $\sqrt T\log T/(8\pi)$;
 \item $R_q$ is an infinite product of positive linear factors, and
 $\Delta_{r,k}(R_q)>0$ for every $r,k\ge1$;
 \item writing $E_q-R_q=\sum_j\varepsilon_{q,j}z^j$, one has
 \begin{equation}
 \sup_{t\ge0}
 \frac{\sum_{j\ge0}|\varepsilon_{q,j}|t^j}{R_q(t)}
 \le\tan^2\!\left(\frac{9\pi}{4q+8}\right)=O(q^{-2}),
 \qquad \Delta_{q,q}(E_q)<0.
 \label{eq:diagonal-control-conclusion}
 \end{equation}
\end{enumerate}
The choices can additionally place every transformed zero in the
usual critical-strip geometry for all sufficiently large $q$.
\end{theorem}

\begin{proof}
Put $m=q-1$, $\theta=9\pi/(4q+8)$, and choose the unique $s>0$ solving
\begin{equation}
 (1+4s)^m(1+2s\cos\theta+s^2)=B.
 \label{eq:diagonal-endpoint}
\end{equation}
The left side increases strictly from one to infinity. Initially use
the finite alphabet
\[
 X=\{se^{i\theta},se^{-i\theta},\underbrace{4s,\ldots,4s}_{m}\},
 \qquad N=q+1,
 \qquad P=\prod_{x\in X}x>0.
\]
Finite elementary coefficient reversal, followed by transposition
and the one-column duality, gives
\begin{equation}
 \det[e_{q+j-i}(X)]_{i,j<q}=P^q h_q(X^{-1}).
 \label{eq:diagonal-reversal}
\end{equation}
Reversal is valid with all out-of-range coefficients zero; repeated
entries introduce no division. After removing the positive factor
$s^{-q}$, it suffices to sign
\[
 H_q=h_q(e^{i\theta},e^{-i\theta},1/4,\ldots,1/4).
\]
Let $Y$ have probability generating function
$((3/4)/(1-w/4))^m$. Its mean and variance are
$\mu=m/3$ and $W=4m/9$. Since the complete coefficients of the unit
pair equal $\sin((j+1)\theta)/\sin\theta$, coefficient extraction gives
\[
 H_q=\frac{(4/3)^m}{\sin\theta}
 \mathbb E\!\left[\sin((q+1-Y)\theta)\mathbf1_{\{Y\le q\}}\right].
\]
Now $(q+1-\mu)\theta=3\pi/2$. The inequality
$\cos x\ge1-x^2/2$ and Chebyshev's inequality imply
\begin{align*}
 \mathbb E[\sin((q+1-Y)\theta)\mathbf1_{\{Y\le q\}}]
 &\le-1+\frac{\theta^2W}{2}+\Pr(Y>q)\\
 &\le-1+\frac{\theta^2W}{2}+\frac{W}{(q-\mu)^2}
 <-1+\frac{13}{q}\le-\frac12.
\end{align*}
Indeed the two positive terms are less than $11.25/q$ and $1/q$,
respectively, using $\pi^2<10$. This makes
\eqref{eq:diagonal-reversal} strictly negative.

To complete the product, let
\[
 Q_J(z)=\prod_{j>J}(1+u_jz),\qquad
 u_j=\left(\frac{\log j}{4\pi j}\right)^2,\qquad J\ge3.
\]
Summability gives $Q_J\to1$ locally uniformly. For each fixed $q$,
choose $J$ sufficiently large that $Q_J(1)<B$, and solve
\eqref{eq:diagonal-endpoint} with $B/Q_J(1)$ in place of $B$, obtaining
$s_J\to s$. Set
\[
 E_q(z)=(1+4s_Jz)^m
 (1+2s_J\cos\theta\,z+s_J^2z^2)Q_J(z).
\]
This retains both endpoints. Its finite coefficient block converges
to the previously constructed block, so its $(q,q)$ determinant
remains strictly negative for all sufficiently large $J$.
This chooses an existent tail cutoff; no effective or certified
numerical cutoff is required for the assertion.

The tail zeros have moduli $(4\pi j/\log j)^2$. Inverting this sequence
gives the count $\sqrt T\log T/(8\pi)$. The convergence exponent is
$1/2$; also $u_j=O(j^{-2+\epsilon})$ for every $\epsilon>0$, so a
split-product bound gives order at most $1/(2-\epsilon)$. The count
gives the reverse order inequality. Thus the order is exactly $1/2$,
and summability gives genus zero. Finite factors change neither the
order nor the count asymptotic.

Replace the one conjugate pair by $(1+yz)^2$, with
$y=\sqrt{1+2s_J\cos\theta+s_J^2}-1$, and call the result $R_q$.
Its infinite positive alphabet gives all the strict reference margins.
The pair agrees at zero and one, and $y<\sqrt B-1<1$.
The exact difference is $d\,z(z-1)T(z)$, where
$d=s_J^2-y^2$ and $T=(1+4s_Jz)^mQ_J$ has positive coefficients.
Consequently
\[
 \frac{\sum_j|\varepsilon_{q,j}|t^j}{R_q(t)}
 \le\frac{dt(t+1)}{(1+yt)^2}
 \le\frac d{y^2}\le\tan^2\theta\qquad(t\ge0),
\]
which proves \eqref{eq:diagonal-control-conclusion}.

Finally $s\sim\log B/(4q)$. Choose $J(q)$ later, if necessary, so
that $s_J/s\to1$ while retaining the strict determinant sign. Then
\[
 \frac{\sin(\theta/2)}{\sqrt{s_J}}\to0,
 \qquad \frac{\cos(\theta/2)}{2\sqrt{s_J}}\to\infty.
\]
The reciprocal pair $\alpha=s_J^{-1}e^{\pm i\theta}$ can therefore
be written as $-(2\rho-1)^2$ with $0<\Re\rho<1$ for all sufficiently
large $q$. Positive real factors map to the critical line. This last
assertion concerns zero geometry only.
\end{proof}

The theorem concerns a family depending on $q$, not one entire function
proved to have infinitely many negative diagonal minors. Scaling both
functions by $c_0=b/B$ supplies the desired values $C(0)=c_0$ and
$C(1)=b$. No nonnegative smooth cosh-transform source, positive Green
primitive, full theta square-source identity, or ordinary Dirichlet
arithmetic is established for the controls. Thus the theorem excludes
a dimension-uniform positivity transfer using only the displayed data.
It does not exclude sharper dimension-dependent conditions, complete
signed compensation, or additional canonical source information.
The separate RGD-E04 r7 assessment accepted all twelve stated
obligations without mathematical correction. This is different-model
internal LLM assessment with shared premises, not specialist review.
